\chapter{Podmartyngały, nadmartyngały i twierdzenia o zbieżności}
\label{ch:nadmartyngaly}

\section{Przewidywalny kierunek zmian a trajektorie}

Z definicji (rozdział~\ref{ch:filtracje})
podmartyngał ma nieujemne warunkowe wartości
oczekiwane przyrostów, a nadmartyngał niedodatnie:
odpowiednio $\E[X_{n+1}-X_n\mid\mathcal F_n]\geq0$
oraz $\E[X_{n+1}-X_n\mid\mathcal F_n]\leq0$ p.n. Nie oznacza to,
że ich trajektorie są odpowiednio niemalejące
lub nierosnące: nawet nadmartyngał może
wiele razy rosnąć, a podmartyngał spadać.
Pojęcia te opisują \emph{warunkową tendencję},
a nie monotoniczność wyników pojedynczego losowania.

\begin{przyklad}[Proces z ujemnym dryfem może rosnąć]
Niech $\xi_k$ będą niezależne, niech
$\mathcal F_n=\sigma(\xi_1,\ldots,\xi_n)$ (z trywialnym
$\mathcal F_0$) i niech $\xi_k$ przyjmują
$+2$ z prawdopodobieństwem $1/4$ oraz $-1$
z prawdopodobieństwem $3/4$. Wtedy
$\E\xi_k=2/4-3/4=-1/4$.
Niech $X_0=0$ i $X_n=\sum_{k=1}^n\xi_k$.
Proces jest adaptowany i całkowalny, gdyż
$|X_n|\leq2n$. Z niezależności
$\E[X_{n+1}-X_n\mid\mathcal F_n]=\E\xi_{n+1}
=-1/4$ p.n., więc $X$ jest nadmartyngałem.
Mimo to w każdym kroku wzrasta o $2$
z prawdopodobieństwem $1/4$.
Proces $M_n=X_n+n/4$ jest martyngałem,
ponieważ jego przyrost ma warunkową średnią zero.
\end{przyklad}

\begin{twierdzenie}[Wypukła funkcja martyngału]
Niech $(M_n)$ będzie rzeczywistym martyngałem
względem $(\mathcal F_n)$,
a $\phi:\R\to\R$ skończoną wypukłą funkcją.
Jeśli $\phi(M_n)\in L^1$ dla każdego $n$,
to $(\phi(M_n))$ jest podmartyngałem
względem tej samej filtracji.
Jeżeli $X$ jest podmartyngałem i $\phi$
jest dodatkowo niemalejąca, to przy
$\phi(X_n)\in L^1$ również $(\phi(X_n))$
jest podmartyngałem.
W szczególności procesy $|M_n|$, $|M_n|^p$
przy odpowiedniej całkowalności ($p\geq1$)
oraz $M_n^2$ dla martyngału w $L^2$
są podmartyngałami.
\end{twierdzenie}
\begin{proof}
Skończona funkcja wypukła na $\R$ jest ciągła,
a więc borelowska, zatem proces $(\phi(M_n))$
jest adaptowany. Całkowalność wynika z założenia.
Warunkowa nierówność Jensena
(rozdział~\ref{ch:condexp}) daje
$\E[\phi(M_{n+1})\mid\mathcal F_n]
\geq\phi(\E[M_{n+1}\mid\mathcal F_n])
=\phi(M_n)$ p.n.
Dla podmartyngału $X$ stosujemy tę samą
nierówność, a następnie monotoniczność:
$\phi(\E[X_{n+1}\mid\mathcal F_n])\geq\phi(X_n)$
p.n. Dla wymienionych funkcji $|x|$, $|x|^p$
($p\geq1$) i $x^2$ stosujemy pierwszą część
twierdzenia z odpowiednim założeniem całkowalności.
\end{proof}

\begin{przyklad}[Wypukłość nie zastępuje całkowalności]
Niech $Y$ ma symetryczny rozkład zadany przez
$\Prob(Y=k)=\Prob(Y=-k)=c/k^3$, $k\geq1$,
gdzie $c=(2\sum_{k=1}^\infty k^{-3})^{-1}$.
Wówczas $\E|Y|=2c\sum_{k\geq1}k^{-2}<\infty$,
$\E Y=0$, lecz $\E Y^2=2c\sum_{k\geq1}k^{-1}=\infty$.
Proces $M_0=0$, $M_n=Y$ dla $n\geq1$, względem
trywialnego $\mathcal F_0$ i $\mathcal F_n=\sigma(Y)$
dla $n\geq1$, jest martyngałem: w chwili zero
warunkowa średnia $Y$ wynosi $\E Y=0$,
a w następnych chwilach proces jest stały.
Jednak
$M_n^2$ ma dla $n\geq1$ nieskończoną wartość oczekiwaną,
a więc nie jest podmartyngałem w definicji
przyjmującej całkowalność. Dlatego przy procesach
$M_n^2$ zakładamy $M_n\in L^2$.
\end{przyklad}

\section{Przekroczenia przedziałów i oscylacje}

Aby dowieść zbieżności prawie na pewno, nie wystarczy
kontrolować średniej procesu. Trzeba wykluczyć
nieskończone oscylacje między dwiema odległymi
wartościami. Służy temu \emph{nierówność
przekroczeń w górę} (upcrossing inequality).

\begin{definicja}[Liczba przekroczeń przedziału]
Niech $x_0,\ldots,x_N\in\R$ oraz $a<b$.
Definiujemy rekurencyjnie
$\sigma_1=\inf\{0\leq k\leq N:x_k\leq a\}$,
$\tau_1=\inf\{\sigma_1<k\leq N:x_k\geq b\}$,
a dla $r\geq2$
$\sigma_r=\inf\{\tau_{r-1}<k\leq N:x_k\leq a\}$
oraz $\tau_r=\inf\{\sigma_r<k\leq N:x_k\geq b\}$,
gdzie indeksy $k$ są całkowite i przyjmujemy
$\inf\varnothing=\infty$.
Po pierwszym nieskończonym czasie wszystkie
następne czasy są nieskończone.
Liczbą ukończonych przekroczeń w górę
od $a$ do $b$ do chwili $N$ nazywamy
$U_N(a,b)=\#\{r:\tau_r\leq N\}$.
Dla procesu $X$ stosujemy tę definicję
ścieżka po ścieżce. Dla ciągu nieskończonego
można równoważnie wyznaczać czasy na całym
$\N_0$ i zliczać tylko te $\tau_r$, które
nie przekraczają $N$: decyzje do chwili $N$
nie zależą od dalszych wartości.
\end{definicja}

\begin{przyklad}[Przekroczenia ukończone i nieukończone]
Dla ścieżki
$(x_0,\ldots,x_8)=(0,2,1,3,0,1,2,0,1)$
oraz $[a,b]=[1,2]$ mamy
$\sigma_1=0$, $\tau_1=1$,
$\sigma_2=2$, $\tau_2=3$,
$\sigma_3=4$, $\tau_3=6$.
Czwarta próba zaczyna się w chwili $\sigma_4=7$,
ale do chwili $8$ nie jest ukończona, więc
$\tau_4=\infty$ w tym horyzoncie.
Zatem $U_8(1,2)=3$. Zliczamy tylko pełne pary,
a nie wszystkie wzrosty ani przekroczenia
pojedynczego poziomu.
\end{przyklad}

\begin{twierdzenie}[Nierówność przekroczeń Dooba]
Niech $(X_n)_{n\geq0}$ będzie rzeczywistym,
całkowalnym podmartyngałem względem
$(\mathcal F_n)$. Dla liczb $a<b$ i $N\in\N_0$
\[
(b-a)\E U_N(a,b)
\leq\E[(X_N-a)^+]-\E[(X_0-a)^+].
\]
\end{twierdzenie}
\begin{proof}
Połóżmy $Y_k=(X_k-a)^+$, gdzie
$x^+=\max(x,0)$. Ponieważ funkcja
$x\mapsto(x-a)^+$ jest wypukła i niemalejąca,
a $0\leq Y_k\leq |X_k|+|a|$ zapewnia całkowalność,
z poprzedniego twierdzenia wynika, że $Y$ jest
nieujemnym podmartyngałem.

Niech $\sigma_r,\tau_r$ będą chwilami określonymi
dla całej nieskończonej trajektorii. Wykażemy,
że są czasami zatrzymania.
Przyjmując pomocniczo $\tau_0=-1$, dla $r\geq1$
i $n\geq0$ mamy
\begin{align*}
\{\sigma_r\leq n\}&=\bigcup_{j=0}^n
 (\{\tau_{r-1}<j\}\cap\{X_j\leq a\}),\\
\{\tau_r\leq n\}&=\bigcup_{j=0}^n
 (\{\sigma_r<j\}\cap\{X_j\geq b\}).
\end{align*}
W pierwszej równości dla $r=1$ zdarzenie
$\{\tau_0<j\}$ jest całym $\Omega$.
Dla nieujemnego czasu zatrzymania $T$
zdarzenie $\{T<j\}$ należy do $\mathcal F_j$:
gdy $j=0$ jest puste, a gdy $j\geq1$ jest
równe $\{T\leq j-1\}\in\mathcal F_{j-1}$.
Indukcja po $r$, najpierw dla $\sigma_r$,
następnie dla $\tau_r$, dowodzi więc wymaganej
mierzalności. W szczególności $U_N(a,b)$ jest
$\mathcal F_N$-mierzalna, ponieważ
$U_N(a,b)=\sum_{r=1}^{N+1}\chi_{\{\tau_r\leq N\}}$:
czasy ukończenia są ściśle rosnące do pierwszej
wartości nieskończonej, a każde ukończone
przekroczenie zajmuje co najmniej jeden krok.

Przed wykonaniem kroku $k$ określmy
\[
H_k=\sum_{r=1}^{N+1}
 \chi_{\{\sigma_r<k\leq\tau_r\}},
 \qquad 1\leq k\leq N.
\]
Nie więcej niż jeden składnik sumy może być równy
$1$, bo przedziały $(\sigma_r,\tau_r]$ są rozłączne.
Ponadto
$\{\sigma_r<k\}=\{\sigma_r\leq k-1\}\in\mathcal F_{k-1}$
oraz $\{\tau_r\geq k\}=\{\tau_r>k-1\}\in\mathcal F_{k-1}$.
Oznacza to, że $H_k\in\{0,1\}$ jest
$\mathcal F_{k-1}$-mierzalny: decyzję o
uwzględnieniu kroku $k$ podejmujemy \emph{przed}
poznaniem przyrostu $Y_k-Y_{k-1}$.

Na każdym przedziale z $\sigma_r<N$
suma wybranych przyrostów teleskopuje do
$Y_{\tau_r\wedge N}-Y_{\sigma_r}$.
Jeśli $\tau_r\leq N$, mamy
$Y_{\sigma_r}=0$ i $Y_{\tau_r}\geq b-a$.
Co najwyżej jeden rozpoczęty przedział może
spełniać $\sigma_r\leq N<\tau_r$; jego wkład
$Y_N-Y_{\sigma_r}=Y_N\geq0$.
Gdy $\sigma_r=N$, wkład ten jest równy zero.
Jeżeli $N=0$, obie sumy przyrostów są puste
i $U_0(a,b)=0$.
Dlatego ścieżkowo
\[
\sum_{k=1}^N H_k(Y_k-Y_{k-1})
 \geq(b-a)U_N(a,b).
\]
Z drugiej strony proces $Y$ jest podmartyngałem,
a więc dla każdego $k$
\[
\E[(1-H_k)(Y_k-Y_{k-1})]
=\E[(1-H_k)\E[Y_k-Y_{k-1}\mid\mathcal F_{k-1}]]\geq0.
\]
Rozpisujemy sumę całkowitych przyrostów $Y$ na
część wybraną przez $H$ i niewybraną:
\begin{align*}
(b-a)\E U_N(a,b)
&\leq\E\sum_{k=1}^NH_k(Y_k-Y_{k-1})\\
&=\E(Y_N-Y_0)
 -\sum_{k=1}^N\E[(1-H_k)(Y_k-Y_{k-1})]\\
&\leq\E(Y_N-Y_0)
=\E[(X_N-a)^+]-\E[(X_0-a)^+].
\end{align*}
Wszystkie sumy są skończone, a ich składniki
całkowalne; nie przechodziliśmy do nieskończonego
horyzontu bez oszacowania.
\end{proof}


\begin{uwaga}[Od liczby przekroczeń do istnienia granicy]
Jeśli ciąg rzeczywisty ma
$\liminf_nx_n<\limsup_nx_n$, to istnieją
liczby wymierne $a<b$ pomiędzy tymi granicami.
Dla dowolnie dużego czasu znajdziemy wówczas
późniejsze wartości $\leq a$ i jeszcze późniejsze
$\geq b$, czyli kolejny przekroczony przedział.
To prowadzi do nieskończenie wielu przekroczeń.
Dlatego skończoność $U_\infty(a,b)$
\emph{dla wszystkich wymiernych $a<b$ jednocześnie}
wyklucza oscylacje; pozostaje jednak osobne
zadanie wykluczenia granic $\pm\infty$.
\end{uwaga}

\begin{twierdzenie}[Twierdzenie Dooba o zbieżności podmartyngałów]
Niech $(X_n)$ będzie rzeczywistym podmartyngałem
i załóżmy, że
$\sup_{n\geq0}\E[X_n^+]<\infty$.
Wtedy istnieje rzeczywista zmienna losowa
$X_\infty\in L^1(\Prob)$ taka, że
$X_n\to X_\infty$ prawie na pewno. Ponadto
$\sup_n\E|X_n|<\infty$.
W szczególności każdy nieujemny nadmartyngał
zbiega prawie na pewno do skończonej zmiennej
losowej, a każdy martyngał z
$\sup_n\E|M_n|<\infty$ zbiega p.n.
\end{twierdzenie}
\begin{proof}
Ustalmy wymierne $a<b$. Z nierówności przekroczeń
oraz nierówności
$(X_n-a)^+\leq X_n^++|a|$
wynika
\[
\sup_N\E U_N(a,b)
\leq\frac{\sup_N\E X_N^++|a|}{b-a}<\infty.
\]
Ponieważ $U_N(a,b)$ jest niemalejąca względem $N$,
definiujemy $U_\infty(a,b)=\lim_N U_N(a,b)$
jako mierzalną zmienną o wartościach w
$\N_0\cup\{\infty\}$. Z twierdzenia o zbieżności
monotonicznej
$\E U_\infty(a,b)<\infty$, a zatem
$U_\infty(a,b)<\infty$ p.n. Dla przeliczalnie
wielu wymiernych par $a<b$ wybieramy jeden
zbiór pełnej miary, na którym wszystkie te liczby
przekroczeń są skończone.

Jeśli na takiej trajektorii zachodziłoby
$\liminf_n X_n<\limsup_n X_n$, można by
wybrać wymierne $a<b$ pomiędzy obiema granicami.
Trajektoria przekraczałaby przedział $[a,b]$
w górę nieskończenie wiele razy, sprzecznie
z uzyskanym wynikiem. Zatem istnieje granica
$X_\infty$ w rozszerzonej prostej $[-\infty,\infty]$.

Pozostaje wykluczyć wartości nieskończone.
Z lematu Fatou
$\E[X_\infty^+]\leq\liminf_n\E X_n^+<\infty$,
więc $X_\infty\ne+\infty$ p.n.
Ponieważ $\E X_n$ jest niemalejące,
$\E X_n\geq\E X_0$ i
$\E X_n^-=\E X_n^+-\E X_n
\leq\sup_k\E X_k^+-\E X_0$.
Dodając ograniczenia na oba momenty części
dodatniej i ujemnej, otrzymujemy
$\sup_n\E|X_n|\leq2\sup_n\E X_n^+-\E X_0<\infty$.
Lemat Fatou zastosowany do $X_n^-$ wyklucza
$X_\infty=-\infty$ z dodatnim prawdopodobieństwem.
To daje skończoną granicę prawie na pewno.
Na mierzalnym zbiorze wyjątkowym definiujemy ją jako zero.
Ponadto oba oszacowania Fatou pokazują, że
$\E X_\infty^+ + \E X_\infty^-<\infty$,
a więc $X_\infty\in L^1(\Prob)$.

Dla nieujemnego nadmartyngału $Z_n$ proces
$X_n=-Z_n$ jest podmartyngałem oraz
$X_n^+=0$. Dla martyngału ograniczonego w $L^1$
mamy $\sup_n\E M_n^+<\infty$.
Oba przypadki są więc konsekwencjami już
udowodnionej wersji twierdzenia.
\end{proof}

\begin{przyklad}[Przyrosty mogą oscylować, mimo zbieżności]
Niech $Z_k$ będą niezależnymi symetrycznymi znakami,
a $M_0=0$, $M_n=\sum_{k=1}^n2^{-k}Z_k$.
W filtracji $\mathcal F_n=\sigma(Z_1,\ldots,Z_n)$
(z trywialnym $\mathcal F_0$) proces jest
adaptowany i ograniczony przez $1$.
Z niezależności warunkowa średnia każdego
przyrostu jest zerowa, więc jest to martyngał.
Szereg
$\sum_k2^{-k}Z_k$ zbiega nawet bezwzględnie
na każdej trajektorii. Zdarzenia $\{Z_{2k-1}=1,Z_{2k}=-1\}$ są niezależne
i mają prawdopodobieństwo $1/4$. Z drugiego lematu
Borela--Cantellego zachodzą nieskończenie często p.n.,
więc dodatnie i ujemne przyrosty występują p.n.
nieskończenie wiele razy.
Zbieżność martyngału nie wymaga monotoniczności.
\end{przyklad}

\section{Jednostajna całkowalność i granice w \texorpdfstring{$L^1$}{L1}}

\begin{definicja}[Jednostajna całkowalność rodziny]
Rodzina $\mathcal K\subseteq L^1(\Prob)$ jest
\emph{jednostajnie całkowalna}, jeśli
\[
\lim_{K\to\infty}\sup_{X\in\mathcal K}
\E[|X|\chi_{\{|X|>K\}}]=0.
\]
\end{definicja}

\begin{lemat}[Rodzina warunkowych wartości jednego $L^1$]
Dla $Y\in L^1(\Prob)$ rodzina
$\{\E[Y\mid\mathcal G]:\mathcal G\subseteq\mathcal F
\text{ jest }\sigma\text{-ciałem}\}$
jest jednostajnie całkowalna.
\end{lemat}
\begin{proof}
Ustalmy $K>0$, $\mathcal G\subseteq\mathcal F$
i połóżmy $Z=\E[Y\mid\mathcal G]$ oraz
$V=|Y|\chi_{\{|Y|>K/2\}}$.
Rozdzielając $Y$ na część o module nie większym
niż $K/2$ i pozostałą, otrzymujemy z warunkowej
nierówności trójkąta
\[
|Z|\leq\frac K2+\E[V\mid\mathcal G].
\]
Na zdarzeniu $A=\{|Z|>K\}\in\mathcal G$
stąd $\E[V\mid\mathcal G]>K/2$.
Po pomnożeniu przez $\chi_{A}$ i wzięciu średniej
wynika nierówność
$\frac K2\Prob(A)\leq\E[V\chi_{A}]\leq\E V$.
Z drugiej strony, znów na $A$, mamy
\[
\E[|Z|\chi_{A}]
\leq\frac K2\Prob(A)
 +\E[\E[V\mid\mathcal G]\chi_{A}]
\leq2\E V.
\]
Zastosowaliśmy tu definiującą własność
warunkowej wartości oczekiwanej przy zdarzeniu
$A\in\mathcal G$. Po prawej znajduje się
wyrażenie niezależne od wyboru $\mathcal G$.
Ponieważ $Y\in L^1$, zbieżność zmajoryzowana daje
$\E[|Y|\chi_{\{|Y|>K/2\}}]\to0$.
Przechodząc z $K\to\infty$ otrzymujemy dokładnie
definicję jednostajnej całkowalności.
\end{proof}

\begin{twierdzenie}[Kryterium $L^1$ dla martyngałów]
Niech $(M_n)$ będzie martyngałem względem
$(\mathcal F_n)$. Następujące warunki są równoważne:
\begin{enumerate}
\item rodzina $\{M_n:n\geq0\}$ jest jednostajnie całkowalna;
\item istnieje $M_\infty\in L^1$ takie, że
$M_n\to M_\infty$ prawie na pewno oraz w $L^1$;
\item istnieje $Y\in L^1$ takie, że
$M_n=\E[Y\mid\mathcal F_n]$ dla każdego $n$.
\end{enumerate}
W warunku (3) można wybrać $Y=M_\infty$.
\end{twierdzenie}
\begin{proof}
Załóżmy (1). Jednostajna całkowalność implikuje
ograniczoność rodziny w $L^1$:
dla dowolnego ustalonego $K$ mamy
$\E|M_n|\leq K+\E[|M_n|\chi_{\{|M_n|>K\}}]$.
Twierdzenie o zbieżności martyngałów daje więc
$M_n\to M_\infty$ p.n. i $M_\infty\in L^1$
na mocy Fatou. Z kryterium Vitaliego (rozdział~\ref{ch:convergence};
zbieżność według prawdopodobieństwa plus
jednostajna całkowalność) otrzymujemy zbieżność
w $L^1$. To daje (2).

Niech zachodzi (2). Dla $m\geq n$ mamy
$M_n=\E[M_m\mid\mathcal F_n]$.
Ponieważ warunkowanie jest kontrakcją w $L^1$,
przechodząc do granicy $m\to\infty$ dostajemy
$M_n=\E[M_\infty\mid\mathcal F_n]$,
czyli (3).

Jeśli zachodzi (3), to poprzedni lemat
zastosowany do tej samej zmiennej $Y$
i do $\mathcal F_n$ daje jednostajną całkowalność
rodziny $(M_n)$, czyli (1).
\end{proof}

\begin{przyklad}[Martyngał o stałej średniej, którego granica ma średnią zero]
Niech $B_1,B_2,\ldots$ będą niezależnymi
próbami Bernoulliego z parametrem $1/2$.
Połóżmy
$M_0=1$ oraz
$M_n=2^n\chi_{\{B_1=\cdots=B_n=1\}}$.
Warunkowo przy historii do chwili $n$
$\E[M_{n+1}\mid\mathcal F_n]=M_n$,
a zatem $\E M_n=1$ dla każdego $n$.
Jednak prawdopodobieństwo, że \emph{wszystkie}
rzuty dały sukces, wynosi
$\lim_n2^{-n}=0$. Zatem $M_n\to0$ prawie na
pewno, choć $\E M_n=1$. Rodzina nie jest
jednostajnie całkowalna, a zbieżność nie
zachodzi w $L^1$. Pojedyncze, coraz rzadsze
i coraz większe wartości niosą całą średnią.
\end{przyklad}

\begin{twierdzenie}[Twierdzenie Lévy'ego o rosnącej informacji]
Niech $(\mathcal F_n)_{n\geq0}$ będzie filtracją
na $(\Omega,\mathcal F,\Prob)$,
$\mathcal F_\infty=\sigma(\bigcup_{n\geq0}\mathcal F_n)$
i $Y\in L^1(\Prob)$. Wtedy
\[
\E[Y\mid\mathcal F_n]
\longrightarrow\E[Y\mid\mathcal F_\infty]
\quad\text{p.n. i w }L^1.
\]
\end{twierdzenie}
\begin{proof}
Niech $M_n=\E[Y\mid\mathcal F_n]$.
Jest to martyngał jednostajnie całkowalny
na mocy poprzedniego lematu; zatem istnieje
$M_\infty\in L^1$ taki, że $M_n\to M_\infty$
p.n. i w $L^1$.
Ponieważ każda $M_n$ jest $\mathcal F_\infty$-
mierzalna, granica ma wersję
$\mathcal F_\infty$-mierzalną.
Dla dowolnego $A\in\mathcal F_k$ i $n\geq k$
\[
\E[M_n\chi_{A}]=\E[Y\chi_{A}].
\]
Przechodząc do granicy w $L^1$, otrzymujemy
$\E[M_\infty\chi_{A}]=\E[Y\chi_{A}]$
dla wszystkich $A$ z $\bigcup_k\mathcal F_k$.
Ta rodzina jest algebrą zbiorów, ponieważ
filtracja jest rosnąca. Klasa zbiorów
$A\in\mathcal F_\infty$ spełniających powyższą
równość jest klasą monotoniczną (wynika to
z zbieżności zmajoryzowanej dla
$M_\infty,Y\in L^1$), więc twierdzenie
o klasie monotonicznej rozszerza równość
na całe $\mathcal F_\infty$.
Z definicji warunkowej wartości oczekiwanej
otrzymujemy
$M_\infty=\E[Y\mid\mathcal F_\infty]$ p.n.
\end{proof}

\begin{przyklad}[Coraz dokładniejsze odczyty obserwacji]
Niech $U\sim\Unif(0,1)$ i
$D_n=2^{-n}\lfloor2^nU\rfloor$,
$\mathcal F_n=\sigma(D_n)$.
Wiedza o $D_n$ oznacza znajomość pierwszych
$n$ cyfr binarnych $U$ (poza nieistotną
kwestią dwóch rozwinięć na zbiorze
prawdopodobieństwa zero). Mamy
$0\leq U-D_n<2^{-n}$, więc $D_n\to U$ p.n.
i $U$ jest $\mathcal F_\infty$-mierzalna.
Twierdzenie Lévy'ego daje
$\E[U\mid\mathcal F_n]\to U$ p.n. i w $L^1$.
Możemy policzyć więcej: przy $D_n=k/2^n$
zmienna $U$ jest jednostajna na przedziale
$[k/2^n,(k+1)/2^n)$, więc
$\E[U\mid\mathcal F_n]=D_n+2^{-(n+1)}$.
Prognoza z obserwacji staje się coraz dokładniejsza.
\end{przyklad}

\section{Zbieżność w \texorpdfstring{$L^2$}{L2} i ograniczona energia}

\begin{twierdzenie}[Martyngał ograniczony w $L^2$]
Niech $(M_n)$ będzie martyngałem z $M_n\in L^2$
i $\sup_n\E M_n^2<\infty$. Wtedy istnieje
$M_\infty\in L^2$ takie, że
$M_n\to M_\infty$ prawie na pewno i w $L^2$.
Ponadto
\[
\sum_{n=1}^{\infty}\E(M_n-M_{n-1})^2
=\E M_\infty^2-\E M_0^2<\infty.
\]
\end{twierdzenie}
\begin{proof}
Dla $k\geq1$ połóżmy $D_k=M_k-M_{k-1}$.
Z ortogonalności przyrostów (rozdział~\ref{ch:martyngaly})
$\sum_{k=1}^n\E D_k^2
=\E M_n^2-\E M_0^2$.
Ponieważ prawa strona jest niemalejąca
i ograniczona, szereg sum kwadratów jest
zbieżny. Dla $m>n$ ortogonalność daje
\[
\E|M_m-M_n|^2=\sum_{k=n+1}^m\E D_k^2\to0
\qquad(n,m\to\infty).
\]
A zatem $(M_n)$ jest ciągiem Cauchy'ego
w zupełnej przestrzeni $L^2$, więc ma granicę
$M_\infty$ w $L^2$. Osobno z nierówności
Cauchy'ego--Schwarza
$\sup_n\E|M_n|\leq\sup_n(\E M_n^2)^{1/2}<\infty$.
Twierdzenie o zbieżności martyngałów daje
zbieżność p.n. do pewnej granicy; granica ta
musi zgadzać się z granicą w $L^2$
(zbieżność w $L^2$ implikuje zbieżność według
prawdopodobieństwa, a granica jest w tym sensie
jednoznaczna). Wreszcie zbieżność w $L^2$
implikuje $\E M_n^2\to\E M_\infty^2$;
przechodzimy do granicy w skończonych
sumach energii.
\end{proof}

\begin{przyklad}[Powrót do procesu rozgałęziającego]
Dla procesu Galtona--Watsona z poprzedniego rozdziału,
przy średniej liczby potomków $m>1$ i
skończonej wariancji $\sigma^2$, martyngał
$W_n=Z_n/m^n$ spełnia
$\sup_n\E W_n^2\leq1+\sigma^2/[m(m-1)]$.
Istnieje więc $W_\infty\in L^2$ takie, że
$W_n\to W_\infty$ p.n. i w $L^2$, a
$\E W_\infty=1$. Nie można jednak na podstawie
samej tej równości stwierdzić, że
$W_\infty>0$ na każdej trajektorii,
na której populacja nigdy nie wymiera;
to wymaga osobnej analizy procesu.
\end{przyklad}

\begin{twierdzenie}[Zbieżność martyngału ograniczonego w $L^p$]
Niech $p>1$ i niech $(M_n)$ będzie martyngałem
spełniającym $\sup_n\E|M_n|^p<\infty$.
Wtedy istnieje $M_\infty\in L^p$ takie, że
$M_n\to M_\infty$ prawie na pewno i w $L^p$.
\end{twierdzenie}
\begin{proof}
Z nierówności Höldera
$\sup_n\E|M_n|\leq(\sup_n\E|M_n|^p)^{1/p}<\infty$.
Twierdzenie Dooba daje więc istnienie skończonej
granicy $M_\infty$ prawie na pewno.
Nierówność maksymalna Dooba w $L^p$ dla
podmartyngału $|M_n|$ mówi, że dla każdego $N$
\[
\E\left(\max_{0\leq k\leq N}|M_k|\right)^p
\leq\left(\frac{p}{p-1}\right)^p\E|M_N|^p.
\]
Prawa strona jest jednostajnie ograniczona.
Z monotonicznej zbieżności przy $N\to\infty$
$M^*:=\sup_{k\geq0}|M_k|$ należy do $L^p$.
Ponieważ $|M_n-M_\infty|^p\leq (2M^*)^p$
prawie na pewno i $M_n\to M_\infty$ p.n.,
zbieżność zmajoryzowana daje
$\E|M_n-M_\infty|^p\to0$.
\end{proof}

\begin{przyklad}[Ograniczenie średniego maksimum]
Jeżeli martyngał ma $\sup_n\E|M_n|^p\leq C$
dla pewnego $p>1$, to jego całkowite
maksimum $\sup_n|M_n|$ ma skończony $p$-ty moment,
a dokładniej
$\E(\sup_n|M_n|)^p\leq(p/(p-1))^pC$.
W przeciwieństwie do samej zbieżności p.n.
otrzymujemy tu kontrolę dużych wychyleń
\emph{całej nieskończonej trajektorii}.
\end{przyklad}

\section{Opcjonalne zatrzymanie przy jednostajnej całkowalności}

\begin{twierdzenie}[Zatrzymanie martyngału domkniętego]
Niech $M_n=\E[Y\mid\mathcal F_n]$, gdzie
$Y\in L^1(\Prob)$, a $\tau$ będzie czasem
zatrzymania z $\Prob(\tau<\infty)=1$.
Wówczas $M_\tau\in L^1$,
$M_\tau=\E[Y\mid\mathcal F_\tau]$ p.n.
i $\E M_\tau=\E M_0$.
\end{twierdzenie}
\begin{proof}
Niech $\tau_n=\tau\wedge n$. Jest to ograniczony
czas zatrzymania, a wzór warunkowania w takim
czasie daje
$M_{\tau_n}=\E[Y\mid\mathcal F_{\tau_n}]$.
Rodzina wszystkich warunkowych wartości tej samej
całkowalnej zmiennej $Y$ jest jednostajnie
całkowalna, zatem $\{M_{\tau_n}:n\geq0\}$
też jest jednostajnie całkowalna.
Ponieważ $\tau<\infty$ p.n., dla niemal każdej
ścieżki od pewnego $n$ mamy $\tau_n=\tau$,
a więc $M_{\tau_n}=M_\tau$. Stąd przez kryterium
Vitaliego $M_{\tau_n}\to M_\tau$ w $L^1$.

Wskażmy dokładnie, dlaczego granica jest warunkową
wartością oczekiwaną. Dla $A\in\mathcal F_\tau$
i każdego $n$ zdarzenie
$A_n=A\cap\{\tau\leq n\}$ należy do
$\mathcal F_{\tau\wedge n}$:
na $\{\tau\wedge n\leq k\}$ przy $k<n$ mamy
$A_n\cap\{\tau\wedge n\leq k\}=A\cap\{\tau\leq k\}
\in\mathcal F_k$, a dla $k\geq n$
$A_n\in\mathcal F_n\subseteq\mathcal F_k$.
Dlatego z definicji warunkowania
\[
 \E[M_{\tau_n}\chi_{A_n}]=\E[Y\chi_{A_n}].
\]
Ponieważ $\tau<\infty$ p.n., mamy
$\chi_{A_n}\to\chi_{A}$ prawie na pewno.
Lewą stronę możemy przeprowadzić do granicy
dzięki zbieżności $M_{\tau_n}\to M_\tau$ w $L^1$
i oszacowaniu
$|\E[M_\tau(\chi_{A_n}-\chi_{A})]|
\leq\E[|M_\tau|\chi_{\{\tau>n\}}]\to0$.
Prawa strona zbiega przez zbieżność zmajoryzowaną,
gdyż $Y\in L^1$. Ostatecznie
$\E[M_\tau\chi_{A}]=\E[Y\chi_{A}]$
dla każdego $A\in\mathcal F_\tau$.
Z osobnej własności mierzalności zatrzymanej
zmiennej $M_\tau$ wynika jej $\mathcal F_\tau$-
mierzalność. Zatem
$M_\tau=\E[Y\mid\mathcal F_\tau]$ p.n., a
w szczególności $\E M_\tau=\E Y=\E M_0$.
\end{proof}

\section{Przykłady przekrojowe i typowe pułapki}


\begin{lemat}[Maksimum nieujemnego nadmartyngału]
Niech $(Z_k)_{0\leq k\leq N}$ będzie
nieujemnym nadmartyngałem względem $(\mathcal F_k)$,
a $a>0$. Wtedy
\[
\Prob\left(\max_{0\leq k\leq N}Z_k\geq a\right)
\leq\frac{\E Z_0}{a}.
\]
\end{lemat}
\begin{proof}
Niech $\tau$ będzie pierwszą chwilą osiągnięcia
poziomu $a$, jeśli występuje przed lub w chwili $N$,
a w przeciwnym wypadku niech $\tau=N$.
Czas $\tau$ jest ograniczonym czasem zatrzymania.
Dla zdarzenia $A=\{\max_{k\leq N}Z_k\geq a\}$
mamy $Z_\tau\geq a$ na $A$, natomiast poza $A$
zmienna $Z_\tau$ pozostaje nieujemna.
Stąd $a\Prob(A)\leq\E Z_\tau$.
Z twierdzenia o opcjonalnym zatrzymaniu
nadmartyngału dla czasu ograniczonego
$\E Z_\tau\leq\E Z_0$, co kończy dowód.
\end{proof}

\begin{przyklad}[Ryzyko przekroczenia progu przed czasem $N$]
Dla martyngału o przyrostach $D_k=M_k-M_{k-1}$
ograniczonych przez deterministyczne stałe $|D_k|\leq c_k$,
$c_k\geq0$, można uzyskać silniejsze
oszacowanie \emph{maksimum} niż przez
sumowanie nierówności dla każdej chwili.
Dla $t>0$ z dowodu Azumy wynika, że
\[
Z_k=\exp\left(t(M_k-M_0)-\frac{t^2}{2}
\sum_{j=1}^k c_j^2\right)
\]
jest nieujemnym nadmartyngałem (warunkowa średnia
czynnika jest co najwyżej jeden).
Jeżeli $\max_{k\leq N}(M_k-M_0)\geq u$, to dla
pewnego $k\leq N$ zachodzi
$Z_k\geq\exp(tu-t^2V_N/2)$, gdzie
$V_N=\sum_{j=1}^N c_j^2$.
Dla nieujemnego nadmartyngału $(Z_k)$
możemy zastosować nierówność maksymalną:
$\Prob(\max_{k\leq N}Z_k\geq a)\leq\E Z_0/a=1/a$.
Wynika ona z ograniczonego zatrzymania przy
pierwszym przekroczeniu $a$ i nieujemności
$Z_N$.
Dla $u>0$ i $V_N>0$ minimalizujemy wykładnik
$-tu+t^2V_N/2$ po $t>0$: minimum przy $t=u/V_N$ daje
\[
\Prob\left(\max_{k\leq N}(M_k-M_0)\geq u\right)
\leq e^{-u^2/(2V_N)}.
\]
Jeżeli $V_N=0$, wszystkie przyrosty do chwili $N$
są zerowe p.n. i dla $u>0$ prawdopodobieństwo po lewej
stronie wynosi zero (bez dzielenia przez $V_N$).
Oszacowanie dotyczy \emph{całej trajektorii do $N$},
a nie tylko wartości końcowej.
\end{przyklad}

\begin{przyklad}[Dlaczego dla nadmartyngału nie trzeba oczekiwać spadku na każdej ścieżce]
Niech $S_n$ będzie symetrycznym spacerem losowym.
Proces $Z_n=S_n^2$ jest podmartyngałem i
$\E Z_n=n$ rośnie, lecz z prawdopodobieństwem
dodatnim $Z_{n+1}<Z_n$ (np. przejście z
$S_n=2$ do $S_{n+1}=1$). Analogicznie
$-S_n^2$ jest nadmartyngałem, choć
może zwiększyć swoją wartość w jednym kroku.
\end{przyklad}

\begin{przyklad}[Inne założenia, inne wnioski]
Trzy fakty trzeba wyraźnie odróżnić:
\begin{enumerate}
\item nieujemny nadmartyngał zbiega p.n.,
  lecz bez dodatkowej kontroli nie musi zbiegać w $L^1$;
\item martyngał ograniczony w $L^1$ zbiega p.n.,
  lecz może stracić wartość oczekiwaną w granicy;
\item martyngał jednostajnie całkowalny zbiega
  p.n. i w $L^1$, a ograniczony w $L^2$
  zbiega nawet w $L^2$.
\end{enumerate}
Przykład $M_n=2^n\chi_{\{B_1=\cdots=B_n=1\}}$
pokazuje jednocześnie, dlaczego nie wolno
zamieniać punktów (2) i (3).
\end{przyklad}

\section{Jak te pojęcia łączą się z łańcuchami Markowa?}

\begin{twierdzenie}[Funkcja harmoniczna łańcucha daje martyngał]
Niech $(X_n)$ będzie jednorodnym łańcuchem Markowa
na przeliczalnym $S$ o macierzy $P$ i naturalnej
filtracji. Niech $h:S\to\R$ będzie ograniczona
i spełnia $Ph=h$, tzn.
$h(i)=\sum_jp_{ij}h(j)$ dla każdego $i\in S$.
Wtedy $(h(X_n))$ jest martyngałem.
Jeżeli $Ph\leq h$, jest nadmartyngałem;
gdy $Ph\geq h$, podmartyngałem.
\end{twierdzenie}
\begin{proof}
Ograniczoność $h$ zapewnia całkowalność.
Z własności Markowa
$\E[h(X_{n+1})\mid\mathcal F_n^X]
=(Ph)(X_n)$.
Podstawienie równości lub odpowiedniej nierówności
prowadzi bezpośrednio do trzech przypadków.
\end{proof}

\begin{przyklad}[Ruina gracza po raz drugi: metoda martyngałowa]
Rozważmy spacer losowy na $\{0,\ldots,N\}$
z granicami pochłaniającymi, z prawdopodobieństwem
kroku w górę $p\in(0,1)$ i w dół $q=1-p$.
Niech $\tau$ będzie chwilą pochłonięcia.
Wykazana wcześniej geometryczna kontrola ogona
$\tau$ daje $\E\tau<\infty$.
Przy $p=q=1/2$ zatrzymany proces $X_{n\wedge\tau}$
jest ograniczonym martyngałem. Zatem przez
zbieżność zmajoryzowaną
$\E_iX_\tau=i$. Ponieważ $X_\tau\in\{0,N\}$,
$N\Prob_i(X_\tau=N)=i$, czyli $h_i=i/N$.
Aby obliczyć średni czas, zauważmy, że
$X_{n\wedge\tau}^2-(n\wedge\tau)$ jest
martyngałem: na $\{\tau>n\}$ mamy
$\E[X_{n+1}-X_n\mid\mathcal F_n]=0$ oraz
$\E[(X_{n+1}-X_n)^2\mid\mathcal F_n]=1$.
Po rozwinięciu różnicy kwadratów jej warunkowa
średnia wynosi zatem $1$ na $\{\tau>n\}$ i $0$
na $\{\tau\leq n\}$, dokładnie tyle samo co przyrost
$(n+1)\wedge\tau-n\wedge\tau$.
Biorąc wartości oczekiwane i przechodząc z
$n\to\infty$, dostajemy
$\E_i X_\tau^2-\E_i\tau=i^2$;
przejście uzasadnia ograniczoność $X_{n\wedge\tau}^2$
i monotoniczna zbieżność $n\wedge\tau\uparrow\tau$.
Ponieważ $\E_iX_\tau^2=N^2(i/N)=Ni$,
$\E_i\tau=i(N-i)$.

Gdy $p\ne q$, proces
$X_{n\wedge\tau}-(p-q)(n\wedge\tau)$ jest
martyngałem. Otrzymujemy
$N h_i-i=(p-q)\E_i\tau$, zgodnie z wcześniejszym
wzorem. Z kolei funkcja harmoniczna
$h_i=[1-(q/p)^i]/[1-(q/p)^N]$
prowadzi do ograniczonego martyngału $h(X_{n\wedge\tau})$
i pozwala odzyskać prawdopodobieństwo wygranej.
Te obliczenia pokazują trzy zastosowania: martyngały liniowe
służą do wyznaczania wartości oczekiwanych, kwadratowe
do obliczania średnich czasów, a harmoniczne do
wyznaczania prawdopodobieństw trafienia.
\end{przyklad}



\section{Powtórzenie: pytania kontrolne o zależność i informację}

\begin{przyklad}[Dlaczego dwie historie nie zawsze dają tę samą prognozę?]
W łańcuchu Markowa, gdy znamy $X_n=i$, każda historia
o dodatnim prawdopodobieństwie kończąca się stanem $i$
prowadzi do tej samej prognozy $X_{n+1}$.
Natomiast w mieszaninie niezależnych prób
Bernoulliego z ukrytym parametrem (przykład z
rozdziału~\ref{ch:markov}) dwa różne przebiegi
mogą kończyć się tym samym wynikiem $X_n$, ale
prowadzą do różnych prognoz, bo inaczej informują
o nieobserwowanym parametrze. To rozróżnienie
pomaga sprawdzić, czy wybrana zmienna stanu
rzeczywiście \emph{opisuje dostatecznie dużo informacji}.

Własność Markowa \emph{nie} wymaga, by $X_{n+1}$
była niezależna od $X_n$. Jedynie warunkowa
zależność od czasów wcześniejszych niż $n$
zanika po poznaniu $X_n$. Inaczej jest przy
zbieżności do stacjonarności: zanikanie zależności
od $X_0$ jest twierdzeniem o odległych czasach
oraz odpowiednich założeniach na macierz przejścia.
\end{przyklad}

